% -----------------------------------------------------------------------------
% prior_diagram_full.tex  --  one-pager sketch of the proximal prior pipeline.
%
% Deliverable T1.1 (fig:prior_pipeline) of docs/prior_appendix_plan.md.
%
% VISUAL GRAMMAR (kept consistent so the figure reads structurally):
%   * SQUARE node      = a parameter  (sampled dimension + fixed MLP + noise root)
%   * CIRCLE node      = a data / random variable  (Xt,Ut,Wt,Zt, U^[A]/U^[Y], A, (Y0,Y1))
%   * ROUNDED BOX      = a process step  (Generate, post-processing, Sample A/Y,
%                        aggregate, validity check, output)
%   * dotted grey      = a parameter configures a step
%   * dotted coloured  = a post-processed block feeds a later generation step
%                        (Xt -> f_U,f_W,f_Z in X's colour; Ut -> f_W,f_Z in U's)
%   * solid arrow      = data flow, in the colour of the block it carries
%   * dashed vermilion = confounding control  (Ut -> U^aug -> U^[A], U^[Y])
%
% COLOUR = MEANING. The palette is the Okabe--Ito colourblind-safe set and is the
%   *same registry* as analysis_prior/design.py::BLOCK_COLORS, so a block keeps
%   one colour across this figure and every appendix plot:
%     X blue - U vermilion - W green - Z orange - A purple - Y sky blue.
%   Anything that is not a structural block (header, post-processing, aggregate,
%   validity check, output, the T*K plate) is deliberately ACHROMATIC, so the
%   causal spine is the only thing carrying hue. Keep it that way: introducing a
%   new hue for a process step would make the reader hunt for a variable that
%   does not exist.
%
% NOTATION: post-processed blocks are written with a tilde, matching the appendix
%   convention where Corrupt(.) + standardisation gives \tilde X, \tilde U, ...
%   Per-mechanism column subsets keep their appendix subscripts (\tilde X_U,
%   \tilde W_Y, \tilde Z_A), and U^[A], U^[Y] are the confounding-controlled inputs.
%   Pill-shaped subset nodes explicitly mark the edge-specific column-subsampling
%   step; the masks only drop parent columns and never add graph edges.
%
% LAYOUT (top -> bottom): header -> parameter row -> generation row ->
%   post-processing box (tilde circles + caption) -> {Sample A | confounding | Sample Y}
%   -> aggregate -> validity check -> output.
%
% X lives on T rows and is only *broadcast* across K, so the green T*K plate is
% L-SHAPED: its top-left notch excludes the Generate-X box and \tilde X; U,W,Z,A,Y
% and the K-aggregate sit inside it. The plate top clears the *whole* of the
% Generate-W,Z box (step 4), which is the tallest node on the generation row.
%
% Code of record: proximalfm/prior/_dataset.py::SCMPrior.{get_batch,generate_dataset},
%                 proximalfm/prior/_reg2cls.py (process_features / dominates),
%                 proximalfm/prior/_prior_config.py (flag defaults).
%
% Bare tikzpicture fragment. Requires (all in docs/latex/preamble.tex):
%   tikz + libraries positioning, arrows.meta, calc.
% If it runs wider/taller than the text block, wrap the \input in
%   \resizebox{\textwidth}{!}{...}  or place it on a dedicated figure page.
%
% NOTE: custom styles are prefixed "pp" to avoid clashing with built-in TikZ keys.
% -----------------------------------------------------------------------------

% ---- palette: Okabe--Ito, mirrored from analysis_prior/design.py -------------
\definecolor{ppcX}{HTML}{0072B2}   % blue       -- baseline covariates
\definecolor{ppcU}{HTML}{D55E00}   % vermilion  -- latent confounder
\definecolor{ppcW}{HTML}{009E73}   % green      -- outcome-inducing proxy
\definecolor{ppcZ}{HTML}{E69F00}   % orange     -- treatment-inducing proxy
\definecolor{ppcA}{HTML}{CC79A7}   % purple     -- treatment
\definecolor{ppcY}{HTML}{56B4E9}   % sky blue   -- outcome

\begin{tikzpicture}[
  line cap=round, line join=round,
  font=\footnotesize,
  % ---- colour-by-block helpers ----------------------------------------------
  % One knob per block: node outlines sit at !70!black (legible at 0.5pt),
  % fills at !7 (a tint, never a wash). #1 is one of ppcX/ppcU/ppcW/ppcZ/ppcA/ppcY.
  ppcol/.style={draw=#1!70!black, fill=#1!7},
  ppink/.style={draw=#1!70!black},                   % same hue, for arrows
  % ---- process boxes (achromatic: structure, not variables) ------------------
  pphdr/.style={rounded corners=2.5pt, draw=black!35, fill=black!4,
    line width=0.5pt, inner sep=5pt, align=center},
  ppgen/.style={rounded corners=2.5pt, ppcol=#1, fill=white,
    line width=0.6pt, inner sep=4pt, align=center},
  ppgen/.default=ppcX,
  ppabox/.style={rounded corners=2.5pt, ppcol=ppcA, fill=ppcA!6,
    line width=0.7pt, inner sep=5pt, align=left, font=\scriptsize},
  ppybox/.style={rounded corners=2.5pt, ppcol=ppcY, fill=ppcY!8,
    line width=0.7pt, inner sep=5pt, align=left, font=\scriptsize},
  ppagg/.style={rounded corners=2.5pt, draw=black!45, fill=white,
    line width=0.6pt, inner sep=5pt, align=center, font=\scriptsize},
  ppacc/.style={rounded corners=2.5pt, draw=black!60, fill=black!5,
    line width=0.7pt, inner sep=5pt, align=center, font=\scriptsize},
  ppout/.style={rounded corners=2.5pt, draw=black!60, fill=white,
    line width=0.7pt, inner sep=5pt, align=center},
  % ---- parameters (squares) -------------------------------------------------
  ppparam/.style={rectangle, ppcol=#1, line width=0.5pt,
    inner sep=3.5pt, font=\scriptsize, minimum height=6mm},
  ppparam/.default=ppcX,
  % ---- data variables (circles) ---------------------------------------------
  ppvar/.style={circle, ppcol=#1, line width=0.7pt,
    inner sep=1pt, minimum size=8mm, font=\footnotesize},
  ppvar/.default=ppcX,
  ppvarU/.style={rounded corners=3pt, ppcol=ppcU, line width=0.7pt,
    inner sep=3pt, font=\scriptsize},
  ppvarY/.style={rounded corners=3pt, ppcol=ppcY, line width=0.7pt,
    inner sep=2.5pt, font=\scriptsize},
  % ---- selected parent-column blocks (compact, coloured pills on data paths) -
  ppsubset/.style={rounded corners=2pt, ppcol=#1, fill=white,
    line width=0.4pt, inner xsep=1.4pt, inner ysep=0.5pt,
    font=\tiny},
  ppsubset/.default=ppcX,
  % ---- containers -----------------------------------------------------------
  ppplate/.style={rounded corners=6pt, draw=black!38, fill=black!3,
    line width=0.7pt},
  ppband/.style={rounded corners=4pt, draw=black!45, fill=none,
    line width=0.6pt, dash pattern=on 1.2pt off 1.6pt},
  % ---- arrows ---------------------------------------------------------------
  ppparr/.style={draw=black!50, line width=0.5pt, dash pattern=on 0.4pt off 1.6pt,
    -{Latex[length=1.5mm,width=1.1mm]}},               % parameter -> step
  ppflow/.style={ppink=#1, line width=0.6pt, -{Latex[length=1.7mm,width=1.2mm]}},
  ppflow/.default=ppcX,
  ppgrey/.style={draw=black!55, line width=0.6pt, -{Latex[length=1.7mm,width=1.2mm]}},
  pparrA/.style={ppink=ppcA, line width=0.6pt, -{Latex[length=1.4mm,width=1mm]}},
  pparrY/.style={ppink=ppcY, line width=0.6pt, -{Latex[length=1.4mm,width=1mm]}},
  ppdup/.style={ppink=ppcU, line width=0.6pt, densely dashed,
    -{Latex[length=1.5mm,width=1.1mm]}},
  ppdep/.style={ppink=#1, line width=0.5pt, densely dotted,
    -{Latex[length=1.4mm,width=1mm]}},                 % gen-step data dependency
  ppdep/.default=ppcX,
  % ---- a light note attached to a variable ----------------------------------
  ppnote/.style={rounded corners=2.5pt, ppcol=ppcU, fill=ppcU!5,
    line width=0.5pt, inner sep=4pt, align=left, font=\scriptsize},
]

% =============================================================================
% Background: the K Monte-Carlo replicate plate (L-shaped; notch excludes the
% X column, which is on T rows and only broadcast across K).
% The top edge sits ABOVE the Generate-W,Z box so step 4 is fully enclosed.
% =============================================================================
\draw[ppplate]
  (0.30,8.20) -- (14.05,8.20) -- (14.05,22.80) -- (4.20,22.80)
  -- (4.20,17.90) -- (0.30,17.90) -- cycle;
\node[font=\scriptsize\itshape, color=black!55, rotate=-90, anchor=south]
  at (14.35,15.0) {computed on the $T\times K$ Monte-Carlo replicate rows};

% =============================================================================
% 1. Header + the parameter squares that configure the generation steps.
% =============================================================================
\node[pphdr, text width=14.8cm] (hdr) at (7.20,25.85) {%
  \textbf{1.\ Sample dimensions, MLP mechanisms, noise roots \& edge masks}\\
  \; {\scriptsize\itshape (once per candidate dataset -- shared across its $K$ replicates)}\\[2pt]
  dimensions $d_X,d_U,d_W,d_Z$ with $\min(d_W,d_Z)\!\ge\!d_U$;\quad
  fixed MLPs $f_X,f_U,f_W,f_Z,f_A,f_Y,f_{\mathrm{aug}}$;\quad
  noise roots $\varepsilon_{(\cdot)}$;\quad edge-specific column masks};

\node[ppparam=ppcX] (pX) at (2.30,24.05) {$d_X,\ f_X,\ \varepsilon_X$};
\node[ppparam=ppcU] (pU) at (7.00,24.05) {$d_U,\ f_U,\ \varepsilon_U$};
\node[ppparam=ppcZ] (pZ) at (10.45,24.05) {$d_Z,\ f_Z,\ \varepsilon_Z$};
\node[ppparam=ppcW] (pW) at (12.85,24.05) {$d_W,\ f_W,\ \varepsilon_W$};

% the (dimension, mechanism, noise) triples are the outputs of step 1.
\draw[ppparr] (hdr.south -| pX) -- (pX.north);
\draw[ppparr] (hdr.south -| pU) -- (pU.north);
\draw[ppparr] (hdr.south -| pW) -- (pW.north);
\draw[ppparr] (hdr.south -| pZ) -- (pZ.north);

% =============================================================================
% 2. Generation row.
% =============================================================================
\node[ppgen=ppcX, text width=2.6cm] (genX)  at (2.30,21.90) {%
  \textbf{2.\ Generate }$X$\\[1pt] from $f_X(\varepsilon_X)$\\[1pt]
  {\scriptsize\itshape on $T$ rows, then broadcast $\times K$}};
\node[ppgen=ppcU, text width=2.9cm] (genU)  at (7.00,21.90) {%
  \textbf{3.\ Generate }$U$\\[1pt] from $f_U(\tilde X_U,\varepsilon_U)$};
\node[ppgen=ppcW, text width=3.4cm] (genWZ) at (11.50,21.90) {%
  \textbf{4.\ Generate }$W,Z$\\[1pt]
  from $f_Z(\tilde U,\tilde X_Z,\varepsilon_Z)$\\ and $f_W(\tilde U,\tilde X_W,\varepsilon_W)$};

% parameter -> step (dotted).
\draw[ppparr] (pX.south) -- (genX.north);
\draw[ppparr] (pU.south) -- (genU.north);
\draw[ppparr] (pZ.south) -- ([xshift=-9mm]genWZ.north);
\draw[ppparr] (pW.south) -- ([xshift= 9mm]genWZ.north);

% =============================================================================
% 3. Post-processing BOX: concise caption on TOP, tilde circles BELOW.
%    Generation arrows terminate on the box's top border; the box maps its four
%    inputs to the post-processed blocks \tilde X, \tilde U, \tilde W, \tilde Z.
%    Transparent fill so the plate notch shows through.
% =============================================================================
\draw[ppband] (0.85,18.20) rectangle (13.75,20.45);
\node[rounded corners=2.5pt, fill=black!5, draw=black!45, line width=0.4pt,
      inner sep=4pt, align=center, font=\scriptsize, text width=11.6cm,
      text=black!75] (pplab) at (7.15,19.75) {%
  \textbf{5.\ Post-process:} categorise $\cdot$ remove outliers $\cdot$ standardise;
  \ \ $W,Z$ categories \emph{dominate} $U$'s};

\node[ppvar=ppcX] (Xt) at (2.30,18.65) {$\tilde X$};
\node[ppvar=ppcU] (Ut) at (7.00,18.65) {$\tilde U$};
% Z sits left of W so that Z -> A (left) and W -> Y (right) do not cross.
\node[ppvar=ppcZ] (Zt) at (10.60,18.65) {$\tilde Z$};
\node[ppvar=ppcW] (Wt) at (12.40,18.65) {$\tilde W$};

% generation -> post-process strip -> post-processed variable.
\draw[ppflow=ppcX] (genX.south)               -- (pplab.north -| genX.south);
\draw[ppflow=ppcX] (pplab.south -| Xt)        -- (Xt.north);
\draw[ppflow=ppcU] (genU.south)               -- (pplab.north -| genU.south);
\draw[ppflow=ppcU] (pplab.south -| Ut)        -- (Ut.north);
\draw[ppflow=ppcZ] ([xshift=-9mm]genWZ.south) -- (pplab.north -| Zt);
\draw[ppflow=ppcZ] (pplab.south -| Zt)        -- (Zt.north);
\draw[ppflow=ppcW] ([xshift= 9mm]genWZ.south) -- (pplab.north -| Wt);
\draw[ppflow=ppcW] (pplab.south -| Wt)        -- (Wt.north);

% -----------------------------------------------------------------------------
% Post-processed blocks feeding the LATER generation steps (dotted dependency):
%   \tilde X -> f_U and f_W,f_Z   (left lane, then a bus under the generation row)
%   \tilde U -> f_W,f_Z           (right lane)
% Both leave their circle ON its border (anchors, not xshifts off the rim).
% -----------------------------------------------------------------------------
\draw[ppdep=ppcX] (Xt.west) -- (0.45,18.65) -- (0.45,20.90) -- (6.05,20.90)
  -- ([xshift=-9.5mm]genU.south)
  node[pos=.79,above,ppsubset=ppcX,xshift=-3mm,yshift=-4mm] {$\tilde X_U$};
\draw[ppdep=ppcX] (6.05,20.90) -- (10.15,20.90) -- ([xshift=-13.5mm]genWZ.south)
  node[pos=.78,above,ppsubset=ppcX,yshift=-4mm] {$\tilde X_Z,\ \tilde X_W$};

\draw[ppdep=ppcU] (Ut.south east) |- (13.95,17.75) -- (13.95,21.90) -- (genWZ.east);

% =============================================================================
% 4. Mechanism steps: Sample A (left), Sample Y0,Y1 (right); confounding control
%    in between.
% =============================================================================
% 6. Confounding control (numbered box) -> U^[A], U^[Y] -> the A and Y mechanisms.
\node[ppnote, text width=4.9cm] (auglab) at (7.00,16.55) {%
  \textbf{6.\ Confounding control.} mix $\tilde U$ with columns of
  $U^{\mathrm{aug}}\!=\!f_{\mathrm{aug}}(\tilde U)$ to reach the sampled shares
  $s_A,s_Y$};
\node[ppvarU] (aug) at (7.00,15.05) {$U^{[A]},\ U^{[Y]}$};

\node[ppabox, text width=4.25cm] (A) at (2.95,12.95) {%
  \textbf{7.\ Sample treatment }$A$\\[2pt]
  $\ell=f_A(U^{[A]},\tilde X_A,\tilde Z_A,\varepsilon_A)$\\[1pt]
  $\tilde\ell=\operatorname{standardise}(\ell)$\\[1pt]
  $A\sim\mathrm{Bern}\!\big(\operatorname{clip}(\sigma(s\tilde\ell+b),\epsilon,1-\epsilon)\big)$};
\node[ppvar=ppcA] (Avar) at (2.95,10.95) {$A$};

\node[ppybox, text width=4.65cm] (Y) at (11.15,12.95) {%
  \textbf{8.\ Sample outcomes }$Y_0,Y_1$\\[2pt]
  $(Y_0^{\mathrm{raw}},Y_1^{\mathrm{raw}})=f_Y(U^{[Y]},\tilde X_Y,\tilde W_Y,\varepsilon_Y)$\\[1pt]
  \emph{heterogeneity remix} with $\gamma,\alpha$};
\node[ppvarY] (Yvar) at (11.15,10.95) {$Y_0,\,Y_1$};

% Parent arrows. The two SHORT ones (Xt -> A, Wt -> Y) stay on their own flank.
% The two LONG ones (Xt -> Y, Zt -> A) are routed orthogonally through a channel
% that passes BELOW the U^[A],U^[Y] node, so nothing crosses on top of it; the two
% channels sit at different heights so they never meet each other either.
% The four X-side departures use angular anchors so they leave \tilde X exactly
% on its rim (an xshift off .south floats free of the circle).
\draw[pparrA] (Xt.250) to[out=-90,in=105] (A.north west)
  node[pos=.52,left,ppsubset=ppcX,yshift=3mm] {$\tilde X_A$};
\draw[pparrY] (Wt.south) to[out=-90,in=75]  (Y.north east)
  node[pos=.52,right,ppsubset=ppcW,yshift=3mm] {$\tilde W_Y$};

\draw[pparrY] (Xt.290) |- (11.60,14.55) -- ([xshift=4.5mm]Y.north)
  node[pos=.44,below,ppsubset=ppcX,yshift=5mm] {$\tilde X_Y$};
\draw[pparrA] (Zt.south) |- (4.95,14.20)  -- ([xshift=20mm]A.north)
  node[pos=.42,below,ppsubset=ppcZ,yshift=5mm] {$\tilde Z_A$};

% confounding control: \tilde U -> box 6 -> U^[A], U^[Y] -> {A,Y}.
\draw[ppdup] (Ut.south) -- (auglab.north)
  node[midway,right,font=\scriptsize,inner sep=1.5pt]{$f_{\mathrm{aug}}$};
\draw[ppdup] (auglab.south) -- (aug.north);
% these leave the node sideways and come down OUTSIDE the parent channels above.
\draw[ppdup] (aug.west) -- (3.60,15.05)  -- ([xshift=6.5mm]A.north);
\draw[ppdup] (aug.east) -- (12.20,15.05) -- ([xshift=10.5mm]Y.north);

% step -> output variable.
\draw[ppflow=ppcA] (A.south) -- (Avar.north);
\draw[ppflow=ppcY] (Y.south) -- (Yvar.north);

% =============================================================================
% 5/6. Aggregate over K + observed-row draw (inside the plate, at the bottom).
% =============================================================================
\node[ppagg, text width=12.2cm] (agg) at (7.10,9.05) {%
  \textbf{9.\ Aggregate over }$K$\textbf{ \& draw the observed row.}\quad
  $\big(\widehat{\mu}_0(x_i),\widehat{\mu}_1(x_i)\big)=\tfrac1K\sum_{k}\big(Y_{0,i,k},Y_{1,i,k}\big)$;\;
  draw $k_i\!\sim\!\mathrm{Unif}\{1,\dots,K\}$, set $y_i=a_i\,y_{1}+(1{-}a_i)\,y_{0}$.};

\draw[ppflow=ppcA] (Avar.south) -- ([xshift=-2.0cm]agg.north);
\draw[ppflow=ppcY] (Yvar.south) -- ([xshift= 2.0cm]agg.north);

% =============================================================================
% 7. Validity check (below the plate).
% =============================================================================
\node[ppacc, text width=13.0cm] (acc) at (7.10,7.45) {%
  \textbf{10.\ Validity check (resample if needed).}\;
  after filtering, require nonempty observed blocks and proxy profiles compatible with $U$.};

\draw[ppgrey] (agg.south) -- (acc.north);

% =============================================================================
% 8. Output (full width, bottom).
% =============================================================================
\node[ppout, text width=13.0cm] (out) at (7.10,5.95) {%
  \textbf{11.\ Output.}\;
  context $\{X,\,W_k,\,Z_k,\,A_k\}$ with observed $Y_k(A_k)$;\quad
  regression targets $\widehat\mu_0,\widehat\mu_1$ (the $K$-averaged potential outcomes from step 9).};

\draw[ppgrey] (acc.south) -- (out.north);

\end{tikzpicture}
